\documentclass[11pt]{article}
\usepackage{ap-lang}
\title{Part 0: The Exam and the Method}
\author{Paul Wang}
\date{November 2026}
\docpart{Part 0 \textperiodcentered\ Study Guide}

\begin{document}
\maketitle

\begin{abstract}
This guide is the operating manual for the AP English Language and Composition exam: what is on it, how each part is scored, and a repeatable method for each question type. It covers the exam's format (Section 1), the four skill categories that organize the course and this series (Section 2), the reading questions (Section 3), the writing questions (Section 4), how to eliminate answer choices (Section 5), the six-point essay rubric shared by all three free-response questions (Section 6), and a study plan (Section 7). Parts 1--4 of the series each take one skill category in depth and end with a full question set and essay. Each section ends with check questions; the key is on the last page.
\end{abstract}

\scoringline{Check questions: \blank\ / \qtotal \qquad Target: 12+ \qquad Weakest section: \blank[10em]}

\section{The exam at a glance}

The exam is fully digital: you read, answer, and type your essays in the Bluebook app. It has two sections.

\begin{compare}{Section}{What it is}{Time and weight}
I. Multiple choice & 45 questions in sets, each set tied to one passage. 23--25 are \textbf{reading} questions on published nonfiction; 20--22 are \textbf{writing} questions on a draft whose sentences are numbered. & 60 minutes; 45\% of the score (about 80 seconds per question). \\
II. Free response & Three essays, each scored out of 6: \textbf{Q1 Synthesis} (argue a position using at least three of six provided sources), \textbf{Q2 Rhetorical analysis} (explain how a writer's choices serve a purpose), \textbf{Q3 Argument} (argue a position using your own evidence). & 2 hours 15 minutes, which includes a 15-minute reading period; 55\% of the score. Plan on about 40 minutes of writing per essay. \\
\end{compare}

Every passage on the exam is nonfiction: speeches, essays, letters, memoirs, journalism, and op-eds from many centuries. You are never asked to identify an author or a historical fact; everything you need is in the passage and its introductory note.

\begin{keyidea}
The three essays are weighted equally, so each one is worth about 18\% of the score, roughly as much as 18 multiple-choice questions. A one-point gain on an essay is a big gain.
\end{keyidea}

\begin{q}{3}{0.1 Format}{Section I: 45 MCQ in 60 minutes (about 80 s each), 45\% of the score.}
Which statement about Section I is accurate?
\ch{It is worth 55\% of the exam score.}{Every question asks you to revise a draft.}{It has 45 questions in 60 minutes, about 80 seconds per question.}{It has 55 questions on fiction and poetry.}
\end{q}

\begin{q}{3}{0.1 Format}{Q1 is the synthesis essay: six sources, use at least three.}
On the free-response section, the essay that supplies six sources and requires you to use at least three of them is the
\chh{rhetorical analysis}{argument essay}{synthesis essay}{reading set}
\end{q}

\section{The four skill categories}

The course is built from nine units, but every unit practices the same four groups of skills, called \textbf{big ideas}. This series takes one per part.

\begin{terms}
Rhetorical situation (Part 1) & Who is writing, to whom, why now, and for what purpose; how the writer's choices fit that situation. \\
Claims and evidence (Part 2) & What the writer asserts, what support is offered, and how commentary links the two. \\
Reasoning and organization (Part 3) & How the claims are sequenced into a line of reasoning; methods of development; transitions. \\
Style (Part 4) & Word choice, comparisons, syntax, and tone, and how they shape the argument's effect. \\
\end{terms}

The multiple-choice questions and the essays test the same skills from two directions: a reading question asks you to \emph{explain} a writer's choice; a writing question asks you to \emph{make} one.

\begin{q}{4}{0.2 Skill categories}{Sequencing claims and transitions = reasoning and organization.}
A question asks why the writer places a personal anecdote before, rather than after, the statistics in paragraph 3. It tests which skill category?
\chh{Style}{Rhetorical situation}{Claims and evidence}{Reasoning and organization}
\end{q}

\section{Reading questions}

\subsection{What they ask}
Reading questions almost never ask what a passage \emph{says}; they ask what a part of it \emph{does}. The common stems:
\begin{terms}
Purpose / function & ``The primary purpose of paragraph 2 is to\ldots''; ``The reference to the shipyard serves mainly to\ldots'' \\
Rhetorical situation & ``The writer's audience is most likely\ldots''; ``The writer acknowledges the audience's\ldots'' \\
Claim and evidence & ``Which claim is supported by the example in\ldots''; ``The statistics function primarily to\ldots'' \\
Organization & ``Which best describes the relationship between paragraphs 3 and 4?'' \\
Style and tone & ``The shift in tone in the final paragraph reflects\ldots''; ``The repetition of `we will' emphasizes\ldots'' \\
Meaning in context & ``As used in paragraph 4, `measured' most nearly means\ldots'' \\
\end{terms}

\subsection{The reading method}
\begin{enumerate}[leftmargin=1.5em]
\item \textbf{Read the introductory note.} It gives the speaker, the audience, the date, and the occasion: half the rhetorical situation, for free.
\item \textbf{Read the passage once for the argument.} Ask: what is the writer's main claim, and what does the writer want this audience to think or do?
\item \textbf{Label each paragraph with its job}, in a few words: ``concedes critics' point,'' ``story of her father,'' ``turns to the solution.''
\item \textbf{Answer function questions in the passage's terms.} Before reading the choices, say what the part does \emph{for the argument}. Then match.
\end{enumerate}

\begin{trap}
The choice that accurately describes the content but not the function. ``The paragraph describes a flooded street'' is true; the answer is what the description does (``makes the abstract risk concrete for readers who have never seen a flood'').
\end{trap}

\begin{q}{2}{0.3 Reading method}{Function questions ask what the part does for the argument, not what it says.}
A question asks for the function of a statistic in paragraph 3. The strongest first step is to
\ch{reread paragraph 3 and list every fact in it}{decide what the statistic does for the writer's claim, then look for the choice that says so}{skip to the last paragraph to find the thesis}{choose the answer that repeats the statistic most precisely}
\end{q}

\begin{q}{4}{0.3 Reading method}{The headnote supplies speaker, audience, occasion, and date.}
The italic note above each reading passage is most useful because it
\ch{summarizes the writer's argument}{is where the questions get their answers}{lists the rhetorical devices in the passage}{supplies the speaker, audience, occasion, and date}
\end{q}

\section{Writing questions}

\subsection{The format}
A writing set shows a draft by a student or other writer, with each sentence numbered: \sn{1}, \sn{2}, and so on. Each question states the writer's goal and asks which revision achieves it. The goal is in the stem; read it twice.
\begin{terms}
Add / keep / delete & ``The writer is considering adding the following sentence after sentence 6. Should the writer add it?'' The choices give yes/no plus a reason; the reason decides. \\
Revise for purpose & ``The writer wants to emphasize the cost to residents. Which version of sentence 9 best does this?'' \\
Evidence & ``Which piece of evidence would best support the claim in sentence 4?'' \\
Placement & ``Where would the following sentence best be placed?'' \\
Transitions & ``Which transition best connects sentences 7 and 8?'' (the logical relationship decides). \\
Thesis and conclusion & ``Which sentence would best serve as the thesis?''; ``Which choice best concludes the passage?'' \\
Audience & ``The writer wants to appeal to an audience of city council members\ldots'' \\
\end{terms}

\subsection{The writing method}
\begin{enumerate}[leftmargin=1.5em]
\item Underline the \textbf{goal} words in the stem (``emphasize,'' ``acknowledge a counterargument,'' ``appeal to parents'').
\item Reread the sentences around the target, at least one before and one after.
\item Choose the option that does \emph{exactly} the stated goal. A beautifully written sentence that does something else is wrong.
\end{enumerate}

\begin{trap}
On ``should the writer add it?'' questions, decide yes or no first, then pick the choice whose \emph{reason} is accurate. Two choices will say ``yes''; usually only one gives the right reason.
\end{trap}

\begin{q}{3}{0.4 Writing questions}{The stated goal decides the answer, not general quality.}
A writing question says the writer wants to ``acknowledge a likely objection from business owners.'' The correct revision will most likely
\ch{add a statistic about the city budget}{restate the thesis more forcefully}{grant that some owners may lose revenue before answering that concern}{use the most vivid language of the four choices}
\end{q}

\begin{q}{1}{0.4 Writing questions}{On add/delete questions the reason must be accurate.}
Two choices say the writer should add a sentence: one because ``it provides an example of the claim in sentence 5,'' the other because ``it introduces the essay's main claim.'' The sentence is an example of sentence 5's claim. Which should you choose?
\ch{The first, because its reason accurately describes the sentence}{Either; both say yes}{Neither; the sentence should be deleted}{The second, because main claims matter most}
\end{q}

\section{Eliminating answer choices}

Wrong choices on this exam follow a few patterns. Name the pattern and the choice is gone.
\begin{terms}
Half right & The first half fits; the second half (``\ldots in order to mock the audience'') does not. One false word sinks the choice. \\
Too strong & ``Proves,'' ``rejects entirely,'' ``all readers,'' ``contemptuous.'' AP writers are usually more measured than the extreme choice. \\
True but irrelevant & Accurate about the passage, but does not answer the question asked (content instead of function). \\
Wrong scope & Describes one sentence when the question asks about the whole paragraph, or vice versa. \\
Reversed & Gets the relationship backward: the example supports the claim, not the claim the example. \\
Outside knowledge & Requires facts about the author or era not in the passage or note. \\
\end{terms}

\begin{keyidea}[Tone choices]
Tone answers are usually two words (``wry and affectionate''). Both words must fit. Extreme negative words (``bitter,'' ``scornful'') are wrong far more often than measured ones (``skeptical,'' ``concerned'').
\end{keyidea}

\begin{q}{2}{0.5 Elimination}{One unsupported half makes a choice wrong.}
A choice reads: ``The writer uses the anecdote to illustrate the danger and to ridicule her opponents.'' The anecdote illustrates the danger, but the passage never mocks anyone. The choice is
\ch{correct, because it names a rhetorical device}{wrong, because it is half right}{wrong, because it is too narrow in scope}{correct, because most of it is accurate}
\end{q}

\begin{q}{1}{0.5 Elimination}{Measured tones are right far more often than extreme ones.}
A passage gently questions a popular policy while praising its goals. Which tone pair is most likely correct?
\chh{measured and skeptical}{gleeful and mocking}{neutral and indifferent}{bitter and scornful}
\end{q}

\section{The essays: one rubric, three rows}

All three essays use the same analytic rubric, out of 6 points.

\begin{compare}{Row}{Earns the point(s)}{Common miss}
A. Thesis (0--1) & A defensible thesis that responds to the prompt: a position someone could reasonably disagree with (Q1, Q3), or an analysis of the writer's choices (Q2). & Restating the prompt; summarizing; a thesis that only names devices (``the writer uses ethos, pathos, and logos''). \\
B. Evidence and commentary (0--4) & 4: specific evidence supports \emph{every} claim in a line of reasoning, and commentary \emph{consistently} explains how the evidence supports it. 3: the same evidence, but commentary explains only some of it. 2: some specific evidence; some explanation, but no clear line of reasoning. 1: general evidence, summarized rather than explained. 0: restates the thesis or the provided material. & Quotes or examples followed by a restatement instead of an explanation of \emph{how} they prove the point. \\
C. Sophistication (0--1) & Sophistication of thought or a complex understanding of the rhetorical situation, shown across the essay: exploring tensions or limits, placing the argument in a broader context, or a consistently vivid, persuasive style. & A single flashy sentence or one sweeping generalization; sophistication must run through the essay. \\
\end{compare}

\subsection{Row B is the essay}
Four of the six points live in Row B, and what moves you up the scale is \textbf{commentary}: the sentences after the evidence that explain why it matters for your claim. A line of reasoning is a chain of claims, each building on the last, that leads to your thesis. Each body paragraph should open with one of those claims, not with a fact.

\begin{keyidea}[The body paragraph pattern]
Claim (a step in your reasoning) $\rightarrow$ specific evidence (a quotation, a source, an example) $\rightarrow$ commentary (how and why the evidence proves the claim; for Q2, how the choice affects the audience and serves the purpose) $\rightarrow$ link back to the thesis.
\end{keyidea}

\subsection{Synthesis-specific rule}
On Q1 the evidence must come from the provided sources: referencing fewer than two of them earns 0 in Row B, and Row B scores of 2--4 require evidence from at least three. Cite each one, e.g.\ (Source C). Your own knowledge may add to the sources but cannot replace them.

\subsection{Timing}
Use the 15-minute reading period for the synthesis sources (annotate which ones support, complicate, or oppose your position). Then budget about 40 minutes per essay: 5 to plan, 32 to write, 3 to reread the thesis and first sentences of paragraphs.

\begin{q}{2}{0.6 Rubric}{Row B is worth 4 of the 6 points.}
Which rubric row carries the most points?
\ch{The rows are equal}{Row B, Evidence and commentary}{Row C, Sophistication}{Row A, Thesis}
\end{q}

\begin{q}{1}{0.6 Rubric}{Summarizing evidence without explaining it caps Row B low.}
An essay quotes the passage often, but after each quotation it restates what the quotation says. Its Row B score is most limited by
\ch{a lack of commentary explaining how the evidence supports the claims}{the absence of a counterargument}{the number of quotations}{the length of the introduction}
\end{q}

\begin{q}{4}{0.6 Rubric}{Q1 Row B scores of 2+ require at least three sources.}
A synthesis essay uses only two of the six sources, though it explains them well. Its Row B score is
\ch{0, because it uses fewer than three sources}{4 if the thesis is defensible}{unaffected, because commentary is what counts}{capped at 1, because scores of 2--4 require evidence from at least three sources}
\end{q}

\section{Study plan}

\begin{compare}{When}{Do this}{With}
Week 1 & Read this guide; take the check questions. Read Part 1 and work its question set untimed. & Parts 0--1 \\
Weeks 2--4 & One part per week: guide first, then the question set timed (about 32 minutes for the multiple choice), then the essay in 40 minutes. & Parts 2--4 \\
Each essay & Score yourself with the part's scoring guide, row by row. Rewrite your weakest body paragraph with better commentary. & Scoring guides \\
Final weeks & Mix sets; take every missed question back to the guide section in its Rule tag. & Answer keys \\
\end{compare}

\begin{q}{1}{0.7 Study plan}{Rewriting with better commentary targets Row B, the 4-point row.}
After scoring an essay at 3/6 (1 thesis, 2 evidence and commentary, 0 sophistication), the most valuable next step is to
\ch{rewrite a body paragraph so the commentary explains how each piece of evidence supports the claim}{memorize more rhetorical device names}{write a longer introduction}{add a quotation from a famous writer}
\end{q}

\makekey
\end{document}
